\documentclass[11pt]{article}

\usepackage[margin=1in]{geometry}
\usepackage{microtype}
\usepackage{amsmath}
\usepackage{graphicx}
\usepackage{booktabs}
\usepackage[hidelinks]{hyperref}

\graphicspath{{figures/}}
\setlength{\emergencystretch}{1.2em}

\newcommand{\natex}{\texttt{natex}}

\title{No Kink at R1: The CVE-Publication Surge Localizes to Late 2025,
Not the DeepSeek-R1 Release}
\author{Hauke Hillebrandt\thanks{University College London. Email:
\texttt{ucjthhi@ucl.ac.uk}. This paper and the underlying \natex{} software were
prepared with substantial assistance from Anthropic's Claude models; the author
reviewed the analyses and text and is responsible for all remaining errors.}}
\date{July 2026 \\ \small analysis pass of record: 2026-07, \natex{} v0.2.0, seed 0}

\begin{document}

\maketitle

\begin{abstract}
\noindent Monthly CVE publications did not bend at DeepSeek-R1. On 90 months
of National Vulnerability Database publication counts (January 2019 -- June
2026), a sharp regression kink design in calendar time at the R1 release (20
January 2025) estimates a slope change in $\ln$ monthly publications of
$+0.038$ ln-points/yr (HC1 se $0.128$, $z=0.30$, $p=0.764$, 95\% CI
$[-0.213, +0.289]$; 18 pre / 17 post months, bandwidth 1.5 years, triangular
kernel). A scout-flagged $+0.18$ ln/yr break --- identical to the
uniform-kernel cell ($+0.184$, se $0.104$, $p=0.077$; Newey--West $p=0.129$
at 6 lags) --- is correctly identified as an artifact of action far from the
cutoff: the estimate grows with bandwidth and vanishes under triangular
weighting; a 49-cutoff placebo grid shows HC1 is oversized on this smooth
serially correlated series (empirical size $0.14$ at nominal $0.05$) and puts
the placebo-calibrated $p$ at $0.80$ ($0.76$ on the 37-cutoff clean subgrid),
with 38 of 49 placebos exceeding the observed $|$kink$|$ and the largest
($+0.385$, $z=3.5$) at November 2020, nowhere near any AI release; a local
quadratic flips the contrast negative ($-0.90$, $p=0.015$ at bandwidth 2.0,
uniform); and sub-window slopes place the first nine post-R1 months at
$+0.103$ ln/yr --- \emph{shallower} than the $+0.280$ pre-trend --- with all
steepening from November--December 2025 onward ($+1.19$ ln/yr through June
2026). Blind LoRD3 localization (seed 0, $k\in\{12,24,36\}$, discovery reads
only the series and time) independently ranks November/December 2025 (scan
$p=0.02$/$0.01$ at $k=24$/$36$) and February--April 2026 first; R1's month
appears in no top-20 at any $k$. An o1-preview contrast is a clean null
($-0.035$, $p=0.844$). The post-2025 acceleration is real --- 3{,}116
publications in November 2025, 5{,}643 in December, 8{,}001 by June 2026 ---
but it arrives 10--17 months after the release, consistent with NVD
operational disruption and backlog catch-up on top of CNA-program expansion.
On this series, attribution of the surge to AI-assisted vulnerability
discovery is unsupported; the result stands as a calibrated null plus an
artifact-identification case study, and reserve-date robustness via the
\texttt{cvelistV5} history is required before any attribution claim.
\end{abstract}

\section{Introduction}\label{sec:intro}

DeepSeek released R1, the first open-weights frontier reasoning model, on 20
January 2025 \cite{deepseek2025}. The release is a natural candidate for a
regime change in the economics of vulnerability discovery: contemporaneous
work showed that LLM agents can autonomously exploit real one-day
vulnerabilities from their CVE descriptions \cite{fang2024}, and an
open-weights reasoning model removes the API chokepoint that had kept such
capabilities meterable. If AI assistance materially accelerated vulnerability
discovery \emph{at scale}, the aggregate flow of published CVEs should
steepen --- kink --- at or shortly after the release. This paper tests that
hypothesis directly.

The series demands discipline before any such reading. CVE publication
counts are heavily process-driven: annual volumes are forecastable to within
roughly 8\% a year ahead from disclosure-pipeline variables alone
\cite{leverett2022}; the publication dates and other fields of the National
Vulnerability Database have documented quality problems \cite{anwar2022};
and NVD processing itself broke and recovered inside the sample window ---
the February 2024 enrichment slowdown and subsequent backlog, followed by
operational changes that NIST announced in April 2026 in response to record
CVE growth \cite{nist2026}. Publication-date slopes therefore conflate
discovery with processing, and any slope break must be localized before it
is attributed.

The estimand --- a change in trend slope at a known date --- is the
regression kink design of \cite{card2015}, applied with calendar time as the
running variable; \cite{bockerman2025} formalize the closely related
difference-in-kinks design that the same toolkit implements. Two
methodological strands discipline the exercise. Regression discontinuity and
kink designs in \emph{time} lack the randomization logic of cross-sectional
designs: serial dependence and smooth trends generate spurious breaks
\cite{hausman2018}. And kink estimates on smooth series are known to produce
spuriously significant slope changes, which is why placebo distributions at
non-event dates are the appropriate yardstick \cite{ganong2018}. Both
concerns are operationalized below --- a 49-cutoff placebo grid whose
calibrated $p$-value is the inference of record, curvature checks, HAC
variants, and a blind localization scan --- estimated with the \natex{}
toolkit \cite{natex}.

The verdict is a clean, well-powered null at the cutoff, plus a correctly
identified artifact: the post-2025 acceleration in CVE publications is real,
but every localization tool places it in late 2025 and 2026, ten to
seventeen months after R1, with the first nine post-release months
\emph{shallower} than the pre-trend.

\section{Data}\label{sec:data}

The outcome series is monthly counts of CVEs by NVD publication date,
obtained from the NVD API 2.0 \cite{nvd} as \texttt{totalResults} per
month-long \texttt{pubStartDate}/\texttt{pubEndDate} window (fetched July
2026): 90 months, January 2019 through June 2026, with no missing months.
The running variable is mid-month decimal time $t = \text{year} +
(\text{month}-0.5)/12$, and the outcome is $\ln$ of the monthly count. The
R1 cutoff is $t=2025.055$ (20 January 2025); the o1-preview contrast cutoff
is $t=2024.70$ (12 September 2024). Levels anchor the story told below:
November 2025 has 3{,}116 publications --- the series' anomalous dip ---
followed by 5{,}643 in December 2025 and 8{,}001 by June 2026. The series
records the NVD \emph{published} date, not the CVE reserve date; the
distinction carries the paper's principal caveat
(Section~\ref{sec:conclusion}).

\section{Design and Methods}\label{sec:methods}

\paragraph{Sharp RKD in calendar time.} The estimand is the change in the
slope of $\ln$ monthly publications, in ln-points per year, at the cutoff,
estimated by local linear regression on each side within the bandwidth
(\natex{} v0.2.0 \cite{natex}, \texttt{natex kink} with a unit policy-kink
change; \cite{card2015,bockerman2025}). The headline specification uses
bandwidth 1.5 years, triangular kernel, degree 1, donut 0 ($n=35$; 18 left,
17 right). A 48-cell grid varies cutoff $\in\{$R1, o1$\}$ $\times$ bandwidth
$\in\{1.0, 1.5, 2.0\}$ years $\times$ kernel
$\in\{\text{triangular},\text{uniform}\}$ $\times$ donut $\in\{0, 1\}$
months $\times$ degree $\in\{1,2\}$. The uniform bandwidth-1.5 cell
reproduces the wave-2 scout's confirmed side-OLS slope delta
($+0.280\rightarrow+0.464$ ln/yr) exactly, and a fully interacted local WLS
mirror of \natex{} matches its uniform kink to $10^{-8}$. The design is
deterministic; seed 0 is declared and no stochastic step exists in
estimation.

\paragraph{Inference on a smooth serially correlated aggregate.} HC1 treats
months as independent; on this series it is demonstrably oversized
(residual lag-1 autocorrelation $0.23$ in the uniform cell), so three guards
are run. \emph{(i) HAC:} Newey--West standard errors (Bartlett, lags 3/6/12)
on the fully interacted local fit, with and without month-of-year dummies.
\emph{(ii) Placebo calibration}, in the spirit of \cite{ganong2018}: the
headline specification re-estimated at 49 monthly placebo cutoffs
(2020.5--2024.5), reporting the empirical size at nominal $0.05$ and the
placebo-calibrated $p = (1 + \#\{|z_{\text{placebo}}| \geq
|z_{\text{R1}}|\})/(1+49)$; a clean 37-cutoff subgrid keeps only placebo
windows that end before R1. The calibrated $p$ is the inference of record.
\emph{(iii) Curvature:} degree-2 refits test whether any kink survives
allowing a smooth quadratic \cite{hausman2018}. \emph{Blind localization:}
\texttt{natex discover} (LoRD3, seed 0, $k\in\{12,24,36\}$) scans the series
with discovery reading only $(t, \ln n)$ --- never the cutoff --- and ranks
candidate break locations; if the R1 date does not surface, the case for an
R1-localized break must come from priors, not the data.

\section{Results}\label{sec:results}

\begin{figure}[t]
\centering
\includegraphics[width=\textwidth]{fig1}
\caption{(a) $\ln$ monthly CVE publications (gray), 2019-01 -- 2026-06, with
per-side local-linear fits at the R1 cutoff (blue, bandwidth 1.5 years,
triangular kernel): the slope change is $+0.038$ ln/yr ($p=0.764$). Orange
dashed segments are sub-window OLS fits: the first nine post-R1 months run
at $+0.103$ ln/yr, \emph{shallower} than the $+0.280$ pre-trend; November
2025 -- June 2026 runs at $+1.19$ ln/yr. Markers along the bottom are the
top blind LoRD3 discovery centers ($k=24$): November/December 2025 and
February--April 2026, plus the April/May 2020 COVID dip --- never R1. The
November 2025 dip (3{,}116) and December spike (5{,}643) are labeled. (b)
The same specification re-estimated at 49 monthly placebo cutoffs
(2020.5--2024.5), with 95\% HC1 intervals; filled markers are nominally
significant at the 5\% level (7 of 49 --- empirical size $0.14$). The R1
estimate (vermillion) sits well inside the placebo distribution: 38 of 49
placebos exceed it in magnitude and the placebo-calibrated $p$ is $0.80$;
the largest placebo ($+0.385$, $z=3.5$, November 2020) is nowhere near any
AI release. The o1 contrast cell (orange) is a clean null.}
\label{fig:fig1}
\end{figure}

\paragraph{Null at the cutoff, with power.} The headline kink at R1
(Table~\ref{tab:kink}, Panel A) is $+0.038$ ln-points/yr (HC1 se $0.128$,
$z=0.30$, $p=0.764$, 95\% CI $[-0.213, +0.289]$): the local trend in $\ln$
publications does not change measurably at the release. The o1-preview
contrast cell --- same specification at $t=2024.70$ --- is likewise null
($-0.035$, se $0.179$, $z=-0.20$, $p=0.844$), matching the scout's confirmed
side-OLS delta of $-0.028$. The null is not for lack of power: the same
specification moved across 49 placebo cutoffs rejects nominally at 7 of 49
and detects a $+0.385$ break ($z=3.5$) at November 2020, so the instrument
registers breaks of the size at issue when they exist.

\paragraph{The uniform $+0.18$ is an edge artifact, and triage identifies
it as one.} The uniform-kernel cell --- numerically identical to the scout's
confirmed side-OLS delta ($+0.280\rightarrow+0.464$) --- gives $+0.184$ (se
$0.104$, $p=0.077$), and is not significant under HAC (Newey--West $p=0.129$
at 6 lags, $0.077$ at 12). Four facts disqualify an R1-localized reading.
First, the kernel/bandwidth pattern: the estimate vanishes under triangular
weighting (which concentrates mass near the cutoff) and \emph{grows} with
bandwidth under uniform weighting ($+0.068$ at bandwidth 1.0, $+0.184$ at
1.5, $+0.217$ at 2.0) --- the signature of action far from the cutoff, not
at it. Second, placebo calibration: HC1's empirical size on this series is
$0.14$ at nominal $0.05$ (7/49; the oversizing catalogued as \natex{} issue
\#51), 38 of 49 placebos exceed the observed $|$kink$|$, and the calibrated
$p$ is $0.80$ on the full grid and $0.76$ on the clean 37-cutoff subgrid.
Third, curvature: degree-2 refits flip the R1 contrast \emph{negative}
($-0.77$ at bandwidth 1.5 triangular; $-0.90$, $p=0.015$, at bandwidth 2.0
uniform) --- a local quadratic absorbs the ``kink'' entirely. Fourth, the
nominally significant stragglers (uniform-with-donut cells at $p=0.011$ --
$0.033$; seasonal Newey--West uniform cells at $p=0.001$ -- $0.033$, the
latter spending 15 parameters on 35 observations) all inherit the same
edge-driven uniform estimate and are overridden by the calibration and the
localization evidence.

\paragraph{Where the action actually is.} Sub-window OLS slopes decompose
the post-period: the pre-R1 window runs at $+0.280$ ln/yr; the first nine
post-R1 months (February--October 2025) run at $+0.103$ --- \emph{shallower}
than the pre-trend; November 2025 -- June 2026 runs at $+1.19$ ln/yr. All
steepening starts roughly ten months after the release. Year-over-year log
growth is $\approx 0.30$ both in the twelve months before R1 and in the last
observed twelve months --- the surge restores the trend as much as it breaks
it. Blind LoRD3 localization agrees: at $k=24$ the top discovery centers are
November/December 2025 (scan $p=0.02$), then February--April 2026; at
$k=36$ the same centers lead (scan $p=0.01$); at $k=12$ the late-2025/2026
cluster leads with the April/May 2020 COVID-era dip next; R1's month appears
in no top-20 at any $k$. The dip-then-surge shape --- November 2025 at
3{,}116 publications (fit residual $-0.45$, the series' anomaly), December
at 5{,}643, June 2026 at 8{,}001 --- is consistent with an NVD operational
disruption followed by backlog catch-up (the October--November 2025 federal
shutdown is the obvious candidate) on top of continuing CNA-program
expansion \cite{nist2026}; that attribution is a hypothesis consistent with
the blind localization, not a result verified in this run.

\begin{table}[t]
\centering
\small
\caption{Sharp RKD in calendar time on $\ln$ monthly CVE publications.
Headline specification: cutoff R1 ($t=2025.055$), bandwidth 1.5 years,
triangular kernel, degree 1, donut 0, HC1.}
\label{tab:kink}
\setlength{\tabcolsep}{4.5pt}
\begin{tabular}{llrrrr}
\toprule
Cell & Kernel/spec & Kink & se & $z$ & $p$ \\
\midrule
\multicolumn{6}{l}{\emph{Panel A: headline cells (ln-points/yr; bandwidth 1.5, degree 1)}}\\
R1 ($t=2025.055$), headline & triangular & $+0.038$ & $0.128$ & $0.30$ & $0.764$ \\
R1, scout's cell            & uniform    & $+0.184$ & $0.104$ & $1.77$ & $0.077$ \\
R1, Newey--West (6 / 12 lags) & uniform  & $+0.184$ & --- & --- & $0.129$ / $0.077$ \\
o1 ($t=2024.70$), contrast  & triangular & $-0.035$ & $0.179$ & $-0.20$ & $0.844$ \\
\midrule
\multicolumn{6}{l}{\emph{Panel B: triage of the uniform estimate}}\\
R1, uniform, bandwidth 1.0 / 2.0 & uniform & $+0.068$ / $+0.217$ & --- & --- & $0.769$ / $0.022$ \\
R1, degree 2, bw 1.5             & triangular & $-0.773$ & $0.529$ & $-1.46$ & $0.144$ \\
R1, degree 2, bw 2.0             & uniform    & $-0.900$ & $0.369$ & $-2.44$ & $0.015$ \\
Largest placebo (Nov 2020)       & triangular & $+0.385$ & $0.109$ & $3.54$ & $0.0004$ \\
\midrule
\multicolumn{6}{l}{\emph{Panel C: placebo calibration of the R1 cell (headline spec)}}\\
Empirical size at nominal $0.05$ & 49 cutoffs & \multicolumn{4}{l}{$0.14$ (7/49; \natex{} issue \#51)} \\
Placebo-calibrated $p$, full / clean & $|z|$ rank & \multicolumn{4}{l}{$0.80$ (49 cutoffs) / $0.76$ (37 clean)} \\
Placebos exceeding observed $|$kink$|$ & --- & \multicolumn{4}{l}{38 of 49} \\
\bottomrule
\end{tabular}

\medskip
\parbox{0.95\textwidth}{\footnotesize Notes: kinks in ln-points per year;
HC1 standard errors against standard normal critical values unless noted.
The placebo grid re-estimates the headline specification at 49 monthly
cutoffs (2020.5--2024.5); the clean subgrid keeps the 37 cutoffs whose
windows end before R1. Sub-window OLS slopes (ln/yr): pre-R1 $+0.280$;
February--October 2025 $+0.103$; November 2025 -- June 2026 $+1.19$.
Blind LoRD3 top centers: November/December 2025 (scan $p=0.02$ at $k=24$,
$0.01$ at $k=36$), then February--April 2026 and April/May 2020; R1 in no
top-20 at any $k$.}
\end{table}

\section{Caveats and Conclusion}\label{sec:conclusion}

Five caveats bound the claims. First, and principally: this is an NVD
\emph{published-date} series. Backlog clearance mechanically shifts counts
late --- the February 2024 NVD crisis and the 2025--26 catch-up are both
in-sample --- so publication-date slopes conflate discovery and processing;
reserve-date robustness via the \texttt{cvelistV5} git history
(\texttt{dateReserved} by CNA, with the Linux-kernel CNA split out) is
mandatory before any attribution claim, and this run did not perform it.
Second, the attribution of the November 2025 dip and December 2025 spike to
the October--November 2025 federal shutdown plus NVD catch-up is a
hypothesis consistent with the blind-scan localization, not verified
in-run. Third, HC1 --- and every nominal $p$-value here --- is oversized on
this smooth serially correlated aggregate (empirical placebo size $0.14$ at
nominal $0.05$; \natex{} issue \#51): the placebo-calibrated inference is
the inference of record, and the few uniform-with-donut and seasonal-HAC
cells at nominal $p=0.001$--$0.033$ inherit the same edge-driven uniform
estimate and are overridden by calibration and localization. Fourth,
calendar time is the running variable: any event co-timed with January 2025
is confounded by construction \cite{hausman2018} --- moot given the null at
the cutoff, but binding for any future positive claim --- and this is a
single national series with no cross-sectional contrast. Fifth, the post
window is truncated at June 2026 (17 post- versus 18 pre-months), and late
placebo-cutoff windows unavoidably overlap the post-R1 region, which is why
the clean 37-cutoff subgrid is reported alongside the full grid.

The conclusion is a calibrated null plus an identified artifact. At the
DeepSeek-R1 release the $\ln$ CVE-publication trend shows no kink
($+0.038$, $p=0.764$; calibrated $p=0.80$) on a design with demonstrated
power; the o1 contrast is equally null. The scout-flagged $+0.18$ is real
arithmetic but wrong attribution: it is the uniform kernel handing the
window edges full weight, and every localization tool --- kernel/bandwidth
pattern, sub-window slopes, curvature absorption, placebo calibration, and
a seed-0 blind scan that never sees the cutoff --- places the action at
November/December 2025 through 2026, ten to seventeen months after the
release, with the first nine post-R1 months \emph{shallower} than the
pre-trend. Attribution of the surge to AI-assisted vulnerability discovery
\cite{fang2024,deepseek2025} is unsupported on this series. The result is
promotable only as an artifact-identification and backlog case study ---
the dip-then-surge is consistent with NVD operational disruption and
catch-up on top of CNA expansion \cite{nist2026} --- and, per the
collection's rules, any attribution-bearing follow-up must start from
reserve dates, not publication dates.

\paragraph{Reproducibility.} All estimates were produced with \natex{}
v0.2.0 \cite{natex}; estimation is deterministic (local weighted least
squares), seed 0 was declared, and the only stochastic step anywhere in the
pipeline is LoRD3's seeded scan calibration. The NVD monthly counts are
fetched by a committed script in the run of record and are not committed to
the repository. Figure~\ref{fig:fig1} regenerates deterministically from the
committed \texttt{figures/make\_fig.py}, which asserts the headline
estimates, the recomputed sub-window slopes, the placebo-calibration
quantities, and the LoRD3 center ranks against the numbers of record before
drawing.

\footnotesize
\begin{thebibliography}{9}
\setlength{\itemsep}{2pt}\setlength{\parskip}{0pt}

\bibitem{anwar2022}
Anwar, A., Abusnaina, A., Chen, S., Li, F., and Mohaisen, D. (2022).
Cleaning the NVD: Comprehensive quality assessment, improvements, and
analyses.
\emph{IEEE Transactions on Dependable and Secure Computing}, 19(6),
4255--4269.

\bibitem{bockerman2025}
B\"ockerman, P., Jysm\"a, S., and Kanninen, O. (2025).
\emph{Difference-in-Kinks Design}.
IZA Discussion Paper No.~18313.
\url{https://docs.iza.org/dp18313.pdf}

\bibitem{card2015}
Card, D., Lee, D.~S., Pei, Z., and Weber, A. (2015).
Inference on causal effects in a generalized regression kink design.
\emph{Econometrica}, 83(6), 2453--2483.

\bibitem{deepseek2025}
DeepSeek-AI (2025).
DeepSeek-R1: Incentivizing reasoning capability in LLMs via reinforcement
learning.
arXiv:2501.12948.

\bibitem{fang2024}
Fang, R., Bindu, R., Gupta, A., and Kang, D. (2024).
LLM agents can autonomously exploit one-day vulnerabilities.
arXiv:2404.08144.

\bibitem{ganong2018}
Ganong, P., and J\"ager, S. (2018).
A permutation test for the regression kink design.
\emph{Journal of the American Statistical Association}, 113(522), 494--504.

\bibitem{hausman2018}
Hausman, C., and Rapson, D.~S. (2018).
Regression discontinuity in time: Considerations for empirical applications.
\emph{Annual Review of Resource Economics}, 10, 533--552.

\bibitem{natex}
Hillebrandt, H. (2026).
\emph{natex: automated natural-experiment discovery and estimation} (version
0.2.0). Software.
\url{https://github.com/HaukeHillebrandt/natex}

\bibitem{leverett2022}
Leverett, \'E., Rhode, M., and Wedgbury, A. (2022).
Vulnerability forecasting: Theory and practice.
\emph{Digital Threats: Research and Practice}, 3(4), Article 42.

\bibitem{nvd}
NIST (2026).
National Vulnerability Database, API 2.0.
\url{https://nvd.nist.gov} (monthly publication counts fetched July 2026).

\bibitem{nist2026}
NIST (2026).
NIST updates NVD operations to address record CVE growth.
News release, April 2026.
\url{https://www.nist.gov/news-events/news/2026/04/nist-updates-nvd-operations-address-record-cve-growth}

\end{thebibliography}

\end{document}
